%% ============================================================
%%  Metadatos del PFG — editar solo esta sección
%% ============================================================
\newcommand{\PFGtitulo}{Evaluación de estrategias de backup ante ransomware mediante laboratorio en contenedores con SIEM Wazuh}
\newcommand{\PFGtitulocorto}{Evaluación impacto ransomware}
\newcommand{\PFGautor}{Alejandro Martinez Martinez}
\newcommand{\PFGdirector}{Antonio Robles}
%\newcommand{\PFGcodirector}{Codirector/a del PFG}
\newcommand{\PFGgrado}{Grado en Ingeniería Informática}
\newcommand{\PFGmodalidad}{Específica}
\newcommand{\PFGcursoacademico}{2026/2027}
\newcommand{\PFGconvocatoria}{Diciembre}
\newcommand{\PFGfechadefensa}{17 de diciembre de 2026}

%% ============================================================
%%  Clase del documento
%% ============================================================
\documentclass[12pt,twoside]{book}

%% ── Codificación y lengua ────────────────────────────────────
\usepackage[utf8]{inputenc}
\usepackage[T1]{fontenc}
\usepackage{lmodern}
\usepackage{xspace}
\usepackage[spanish,es-tabla]{babel}
\spanishplainpercent %% El documento ya inserta \,\% manualmente; evita conflicto de glue con babel-spanish

%% ── Márgenes ─────────────────────────────────────────────────
\usepackage[papersize={210mm,297mm},
            tmargin=25mm, bmargin=25mm,
            lmargin=20mm, rmargin=20mm]{geometry}

%% ── Párrafos ─────────────────────────────────────────────────
\setlength\parindent{0cm}
\setlength\parskip{1em}

%% ── Colores (tabla ANTES de xcolor para que funcione rowcolors) ──
\usepackage[dvipsnames,table]{xcolor}
\definecolor{links}         {rgb}{0.35, 0.4, 0.35}
\definecolor{urls}          {rgb}{0,0,1}
\definecolor{cites}         {rgb}{0.8,0,0}
\definecolor{gray20}        {gray}{0.20}
\definecolor{tabla-cabecera}{RGB}{31,56,100}
\definecolor{tabla-par}     {RGB}{235,240,250}
\definecolor{codegreen}     {rgb}{0,0.6,0}
\definecolor{codegray}      {rgb}{0.5,0.5,0.5}
\definecolor{codepurple}    {rgb}{0.58,0,0.82}
\definecolor{backcolour}    {rgb}{0.95,0.95,0.92}

%% ── Gráficos ─────────────────────────────────────────────────
\usepackage{graphicx}
\usepackage{float}
\usepackage{caption}
\usepackage{subcaption}
\graphicspath{{imagenes/}}

%% ── Hipervínculos ────────────────────────────────────────────
\usepackage[hyphens]{url}
\usepackage{hyperref}
\hypersetup{
    colorlinks = true,
    linkcolor  = links,
    urlcolor   = urls,
    citecolor  = cites,
    pdftitle   = {\PFGtitulo},
    pdfauthor  = {\PFGautor}
}

%% ── Tablas ───────────────────────────────────────────────────
\usepackage{booktabs}
\usepackage{tabularx}
\usepackage{multicol}
\usepackage{multirow}
\usepackage{colortbl}
\usepackage{longtable}
\usepackage{array}
\usepackage{siunitx}
\renewcommand{\arraystretch}{1.5}

%% ── Matemáticas ──────────────────────────────────────────────
\usepackage{amsmath}
\usepackage{amssymb}

%% ── Resaltado ────────────────────────────────────────────────
\usepackage{soul}

%% ── Código fuente ────────────────────────────────────────────
\usepackage{listings}
\lstloadlanguages{C,C++,Python,Java,PHP,HTML,bash}
\lstdefinestyle{mystyle}{
    commentstyle    = \color{codegreen},
    keywordstyle    = \color{blue},
    numberstyle     = \tiny\color{codegray},
    stringstyle     = \color{codepurple},
    backgroundcolor = \color{backcolour},
    basicstyle      = \footnotesize\ttfamily\color{black},
    breaklines      = true,
    captionpos      = b,
    frame           = single,
    keepspaces      = true,
    numbers         = left,
    numbersep       = 10pt,
    showstringspaces= false,
    stepnumber      = 1,
    tabsize         = 2,
    xleftmargin     = 1.5cm,
}
\lstset{style=mystyle}
\renewcommand{\lstlistingname}{Listado}
\renewcommand{\lstlistlistingname}{Índice de listados}
\lstset{
    extendedchars=true,
    inputencoding=utf8,
    literate=
        {á}{{\'a}}1 {é}{{\'e}}1 {í}{{\'i}}1 {ó}{{\'o}}1 {ú}{{\'u}}1
        {Á}{{\'A}}1 {É}{{\'E}}1 {Í}{{\'I}}1 {Ó}{{\'O}}1 {Ú}{{\'U}}1
        {ñ}{{\~n}}1 {Ñ}{{\~N}}1 {ü}{{\"u}}1 {Ü}{{\"U}}1
        {¿}{{¿}}1 {¡}{{¡}}1
}
\lstdefinestyle{bash}{
    language        = bash,
    commentstyle    = \color{codegreen},
    keywordstyle    = \color{blue},
    basicstyle      = \footnotesize\ttfamily\color{black},
    backgroundcolor = \color{backcolour},
    breaklines      = true,
    frame           = single,
    numbers         = left,
    numbersep       = 10pt,
    xleftmargin     = 1.5cm,
    captionpos      = b,
}

%% ── Encabezados (según reglamento UNED) ─────────────────────
%% Pares:   nº pág (izq) | título PFG (der)
%% Impares: título cap (izq) | nº pág (der)
\usepackage{fancyhdr}
\fancyhf{}
\lhead[\thepage]{\leftmark}
\rhead[\PFGtitulocorto]{\thepage}
\setlength{\headheight}{16pt}
\pagestyle{empty}
%% Primera página de capítulo sin encabezado
\fancypagestyle{plain}{
  \fancyhf{}
  \fancyfoot[C]{\small\thepage}
  \renewcommand{\headrulewidth}{0pt}
}

%% ── Formato de capítulos ─────────────────────────────────────
\usepackage{titlesec}
\titleformat{\chapter}[display]
    {\normalfont\huge\bfseries}{\chaptertitlename\ \thechapter}{20pt}{\Huge}
\titlespacing*{\chapter}{0pt}{-50pt}{40pt}
\setcounter{secnumdepth}{4}
\titleformat{\paragraph}
    {\normalfont\normalsize\bfseries}{\theparagraph}{1em}{}
\titlespacing*{\paragraph}
    {0pt}{3.25ex plus 1ex minus .2ex}{1.5ex plus .2ex}

%% ── Bibliografía en índice ───────────────────────────────────
\usepackage[nottoc,notlot,notlof]{tocbibind}

%% ── Glosario y acrónimos ─────────────────────────────────────
\usepackage[acronym,toc]{glossaries}
\makeglossaries
% acronimos.tex
\newacronym{nist}{NIST}{Instituto Nacional de Estándares y Tecnología}
\newacronym{cis}{CIS}{Centro de Seguridad de Internet}
\newacronym{ca}{CA}{Autoridad de Certificación}
\newacronym{edr}{EDR}{Detección y Respuesta en el Punto Final}
\newacronym{siem}{SIEM}{Gestión de Eventos e Información de Seguridad}
\newacronym{fim}{FIM}{Monitorización de la Integridad de Ficheros}
\newacronym{mttr}{MTTR}{Tiempo Medio de Recuperación}
\newacronym{rto}{RTO}{Objetivo de Tiempo de Recuperación}
\newacronym{rpo}{RPO}{Objetivo de Punto de Recuperación}
\newacronym{raas}{RaaS}{Ransomware como Servicio}
\newacronym{apt}{APT}{Amenaza Persistente Avanzada}
\newacronym{aes}{AES}{Estándar de Cifrado Avanzado}
\newacronym{gcm}{GCM}{Modo Galois/Contador}
\newacronym{sftp}{SFTP}{Protocolo de Transferencia de Ficheros sobre SSH}
\newacronym{ssh}{SSH}{Shell Seguro}
\newacronym{s3}{S3}{Servicio de Almacenamiento Simple}
\newacronym{iam}{IAM}{Gestión de Identidades y Accesos}
\newacronym{api}{API}{Interfaz de Programación de Aplicaciones}
\newacronym{json}{JSON}{Notación de Objetos de JavaScript}
\newacronym{yaml}{YAML}{YAML no es un Lenguaje de Marcado}
\newacronym{xml}{XML}{Lenguaje de Marcado Extensible}
\newacronym{csv}{CSV}{Valores Separados por Comas}
\newacronym{cli}{CLI}{Interfaz de Línea de Comandos}
\newacronym{tls}{TLS}{Seguridad de la Capa de Transporte}
\newacronym{smb}{SMB}{Bloque de Mensajes de Servidor}
\newacronym{xdr}{XDR}{Detección y Respuesta Extendidas}
\newacronym{ids}{IDS}{Sistema de Detección de Intrusiones}
\newacronym{ips}{IPS}{Sistema de Prevención de Intrusiones}
\newacronym{vm}{VM}{Máquina Virtual}
\newacronym{os}{OS}{Sistema Operativo}
\newacronym{cpu}{CPU}{Unidad Central de Procesamiento}
\newacronym{ram}{RAM}{Memoria de Acceso Aleatorio}
\newacronym{nsa}{NSA}{Agencia de Seguridad Nacional (EE.\,UU.)}
\newacronym{ddos}{DDoS}{Denegación de Servicio Distribuida}
\newacronym{worm}{WORM}{Escritura Única, Lectura Múltiple}
\newacronym{ueba}{UEBA}{Análisis del Comportamiento de Usuarios y Entidades}
\newacronym{rosi}{ROSI}{Retorno de la Inversión en Seguridad}
\newacronym{hids}{HIDS}{Sistema de Detección de Intrusiones basado en Host}
\newacronym{jvm}{JVM}{Máquina Virtual de Java}
\newacronym{sdn}{SDN}{Redes Definidas por Software}
\newacronym{nat}{NAT}{Traducción de Direcciones de Red}
\newacronym{bore}{BORE}{Potenciador de Respuesta Orientado a Ráfagas}
\newacronym{ects}{ECTS}{Sistema Europeo de Transferencia y Acumulación de Créditos}
\newacronym{rgpd}{RGPD}{Reglamento General de Protección de Datos}
\newacronym{hipaa}{HIPAA}{Ley de Portabilidad y Responsabilidad del Seguro Médico (EE.\,UU.)}


%% ── Anexos ───────────────────────────────────────────────────
\usepackage{appendix}
\renewcommand{\appendixname}{Anexo}
\renewcommand{\appendixtocname}{Anexos}
\renewcommand{\appendixpagename}{Anexos}

%% ── Subfiles ─────────────────────────────────────────────────
\usepackage{subfiles}

%% ── Enumeraciones ────────────────────────────────────────────
\usepackage{enumitem}

%% ── Floats y posicionamiento ─────────────────────────────────
\usepackage{placeins}
\usepackage{needspace}

%% ── Nombres en español ───────────────────────────────────────
\renewcommand{\spanishtablename}{Tabla}
\renewcommand{\spanishlisttablename}{Índice de tablas}
\renewcommand{\spanishcontentsname}{Índice}
\renewcommand{\spanishfigurename}{Figura}
\renewcommand{\spanishlistfigurename}{Índice de figuras}

%% ── Tipo de columna centrada con ancho fijo ──────────────────
\newcolumntype{C}[1]{>{\centering\arraybackslash}p{#1}}

%% ── Comandos auxiliares generales ────────────────────────────
\newcommand\paginablanco{%
    \null\thispagestyle{empty}\newpage}
\newcommand\paginablancosin{%
    \paginablanco\addtocounter{page}{-1}}
\newcommand{\pfg}{Proyecto Fin de Grado\xspace}
\newcommand\ddfrac[2]{\frac{\displaystyle #1}{\displaystyle #2}}

%% ── Comandos del PFG ─────────────────────────────────────────
\newcommand{\restic}{Restic\xspace}
\newcommand{\wazuh}{Wazuh\xspace}
\newcommand{\tdetect}{\ensuremath{T_{\text{detect}}}}
\newcommand{\tlaunch}{\ensuremath{T_{\text{launch}}}}
\newcommand{\mttr}{\gls{mttr}}
\newcommand{\stratA}{Estrategia~A}
\newcommand{\stratB}{Estrategia~B}
\newcommand{\stratC}{Estrategia~C}

%% Compatibilidad \textcite (bibliografía numérica IEEEtran vía bibtex)
\newcommand{\textcite}[1]{\cite{#1}}

%% ============================================================
\begin{document}
%% ============================================================

\pagenumbering{gobble}

%% ── PORTADA (modelo B.9) ─────────────────────────────────────
\begin{titlepage}
\begin{figure}[!htb]
    \centering
    \includegraphics[width=4cm]{rosadelosvientos.png}
\end{figure}
\begin{center}
    {\scshape\Large Universidad Nacional de Educación a Distancia \par}
    {\scshape\Large Escuela Técnica Superior de Ingeniería Informática \par}
    \vfill
    \large{Proyecto de Fin de Grado en \PFGgrado}\\
    \vfill
    {\bfseries\Huge \textcolor{black}{\PFGtitulo} \par}
    \vfill
\end{center}
\vspace*{2cm}
{\huge \textcolor{black}{\PFGautor} \par}
{\huge \textcolor{black}{Dirigido por: \PFGdirector} \par}
\ifdefined\PFGcodirector
    {\huge \textcolor{black}{Codirigido por: \PFGcodirector} \par}
\fi
{\huge \textcolor{black}{Curso \PFGcursoacademico, convocatoria \PFGconvocatoria} \par}
\end{titlepage}

%% ── PRIMERA HOJA (modelo B.10) ───────────────────────────────
\chapter*{}
\thispagestyle{empty}
\begin{figure}[!htb]
    \centering
    \includegraphics[width=4cm]{rosadelosvientos.png}
\end{figure}
\begin{center}
    {\bfseries\Huge \textcolor{black}{\PFGtitulo} \par}
    \vfill
    {\bfseries\Large \textcolor{black}{Proyecto de Fin de Grado en \PFGgrado} \par}
    {\bfseries\Large \textcolor{black}{de modalidad \PFGmodalidad} \par}
\end{center}
\vspace*{4cm}
{\Large Realizado por: \PFGautor \par}
{\Large Dirigido por: \PFGdirector \par}
\ifdefined\PFGcodirector
    {\Large Codirigido por: \PFGcodirector \par}
\fi
\vfill
{\Large Fecha de lectura y defensa: \PFGfechadefensa \par}
\newpage

%% ── PRELIMINARES ─────────────────────────────────────────────
\pagenumbering{Roman}

%% Agradecimientos
\chapter*{Agradecimientos}
\thispagestyle{empty}
\addcontentsline{toc}{chapter}{Agradecimientos}

\begin{flushright}
\textit{A mi familia, por el apoyo incondicional durante estos años.}
\end{flushright}

Este trabajo no habría sido posible sin la orientación de mi director, cuya
disponibilidad y criterio técnico han sido fundamentales en cada etapa del proyecto.
Agradezco también a la comunidad de software libre que desarrolla y mantiene las
herramientas sobre las que se construye este laboratorio: Wazuh, MinIO, Restic y
Docker, entre otras.

\cleardoublepage

%% Resumen
\chapter*{Resumen}
\addcontentsline{toc}{chapter}{Resumen}

El ransomware se ha convertido en una de las amenazas más graves para organizaciones
de todo tipo. Los ataques modernos no se limitan a cifrar los datos de la víctima:
incluyen sistemáticamente la destrucción de las copias de seguridad antes del
cifrado, eliminando la principal vía de recuperación. Frente a este escenario, la
industria propone el almacenamiento inmutable como mecanismo de protección, pero la
literatura académica carece de comparativas experimentales que cuantifiquen su
impacto real en el tiempo de recuperación.

Este Trabajo de Fin de Grado presenta el diseño, implementación y evaluación de un
laboratorio de simulación de ransomware basado íntegramente en contenedores Docker.
El entorno reproduce una infraestructura empresarial mínima compuesta por un
servidor de ficheros, una \textit{workstation}, un servidor de \textit{backups}
SFTP, un sistema de almacenamiento de objetos con Object Lock (MinIO) y un SIEM
completo (Wazuh~4.7.5 con OpenSearch).

El diseño experimental sigue un esquema factorial de cuatro escenarios de ataque
por tres estrategias de recuperación y veinte repeticiones, totalizando
240~ejecuciones completamente automatizadas y ejecutadas en orden aleatorio. Los
escenarios cubren las tácticas más relevantes de la matriz MITRE ATT\&CK: cifrado
masivo de ficheros (T1486), destrucción activa de copias de seguridad (T1490),
movimiento lateral por SSH (T1021.004) y exfiltración de datos (T1041). Las tres
estrategias evaluadas son: ausencia de copia de seguridad, repositorio Restic sobre
SFTP convencional y repositorio Restic sobre MinIO con Object Lock COMPLIANCE.

El hallazgo central es de disponibilidad: ante la destrucción activa del
repositorio, la estrategia SFTP resulta irrecuperable en el 100\,\% de los casos,
mientras que la combinación de Object Lock y una política de mínimo privilegio
mantiene una tasa de éxito del 100\,\% (test exacto de Fisher $p < 0.001$). Además,
cuando se mide de forma aislada la velocidad del \textit{backend}, la restauración
desde el almacenamiento inmutable resulta un 39\,\% más rápida que desde SFTP, con
un tamaño del efecto máximo ($\delta$ de Cliff $=1.0$). Los tiempos de detección
mediante Wazuh, medidos a partir del instante real de la alerta, oscilan entre 0.4
y 11.2~segundos, muy por debajo del umbral de un minuto y en marcado contraste con
los tiempos de permanencia de meses reportados por la industria para incidentes no
instrumentados. El laboratorio, con detección por monitorización de integridad de
ficheros (FIM) y eventos estructurados, reglas de correlación Wazuh \textit{custom}
y recuperación verificada por \textit{hash}, se publica como software libre.

\cleardoublepage

%% Abstract
\chapter*{Abstract}
\addcontentsline{toc}{chapter}{Abstract}

Ransomware has become one of the most serious threats to organizations of all types.
Modern attacks do not merely encrypt the victim's data: they systematically destroy
backup copies before encryption, eliminating the primary recovery path. In response,
the industry proposes immutable storage as a protection mechanism, yet the academic
literature lacks experimental comparisons that quantify its actual impact on recovery
time.

This Bachelor's Thesis presents the design, implementation and evaluation of a
ransomware simulation laboratory based entirely on Docker containers. The environment
replicates a minimal enterprise infrastructure consisting of a file server, a
workstation, an SFTP backup server, an object storage system with Object Lock
protection (MinIO), and a complete SIEM stack (Wazuh~4.7.5 with OpenSearch).

The experimental design follows a factorial scheme of four attack scenarios, three
recovery strategies and twenty repetitions, totalling 240~fully automated runs
executed in randomized order. The scenarios cover the most relevant MITRE ATT\&CK
tactics: mass file encryption (T1486), active destruction of backup copies (T1490),
lateral movement via SSH (T1021.004), and data exfiltration (T1041). The three
strategies evaluated are: no backup, a Restic repository over conventional SFTP, and
a Restic repository over MinIO with COMPLIANCE Object Lock.

The central finding concerns availability: under active destruction of the
repository, the SFTP strategy is unrecoverable in 100\,\% of cases, whereas the
combination of Object Lock and a least-privilege policy maintains a 100\,\% success
rate (Fisher exact test $p < 0.001$). In addition, when the backend speed is
measured in isolation, restoring from immutable storage is 39\,\% faster than from
SFTP, with a maximal effect size (Cliff's $\delta = 1.0$). Detection times via
Wazuh, measured from the actual alert instant, range from 0.4 to 11.2~seconds, well
below the one-minute threshold and in sharp contrast with the months-long dwell
times reported by the industry for non-instrumented incidents. The laboratory,
featuring FIM- and event-based detection, custom Wazuh correlation rules and
hash-verified automated recovery, is published as free software.

\cleardoublepage

%% Palabras clave
\chapter*{Palabras clave}
\addcontentsline{toc}{chapter}{Palabras clave}

ransomware, Object Lock, SIEM, \textit{backup} inmutable, MITRE ATT\&CK.

\cleardoublepage

%% Keywords
\chapter*{Keywords}
\addcontentsline{toc}{chapter}{Keywords}

ransomware, Object Lock, SIEM, immutable backup, MITRE ATT\&CK.

\cleardoublepage

%% ── Índices ──────────────────────────────────────────────────
\renewcommand{\listtablename}{Índice de tablas}
\renewcommand{\listfigurename}{Índice de figuras}

\tableofcontents
\cleardoublepage

\listoftables
\addcontentsline{toc}{chapter}{Índice de tablas}
\cleardoublepage

\listoffigures
\addcontentsline{toc}{chapter}{Índice de figuras}
\cleardoublepage

%% Lista de acrónimos
\printglossary[type=\acronymtype, title={Lista de acrónimos},
               toctitle={Lista de acrónimos}]
\cleardoublepage

\paginablanco{}

%% ── CUERPO DEL DOCUMENTO ─────────────────────────────────────
\pagenumbering{arabic}
\pagestyle{fancy}

\subfile{capitulos/01_introduccion}
\subfile{capitulos/02_estado_arte}
\subfile{capitulos/03_arquitectura}
\subfile{capitulos/04_disenyo_experimental}
\subfile{capitulos/05_resultados}
\subfile{capitulos/06_discusion}
\subfile{capitulos/07_conclusiones}

%% ── BIBLIOGRAFÍA ─────────────────────────────────────────────
\newpage
\bibliography{bibliografia}
\bibliographystyle{IEEEtran}
\clearpage

%% ── ANEXOS ───────────────────────────────────────────────────
\appendix
\subfile{capitulos/anexo_a_reglas}
\subfile{capitulos/anexo_b_art}
\subfile{capitulos/anexo_c_repositorio}

\end{document}
